\documentclass[11pt,a4paper]{article}
\usepackage[T1]{fontenc}
\usepackage[utf8]{inputenc}
\usepackage[margin=23mm]{geometry}
\usepackage{lmodern,microtype}
\usepackage{amsmath,amssymb}
\usepackage{booktabs,longtable,array,tabularx}
\usepackage{enumitem}
\usepackage[hyphens]{url}
\usepackage[hidelinks]{hyperref}
\usepackage{fancyhdr}
\setlength{\headheight}{14pt}
\pagestyle{fancy}
\fancyhf{}
\lhead{Independent academic review for IEEE TIFS --- thirteenth cycle}
\rhead{11 September 2026}
\cfoot{\thepage}
\setlength{\parindent}{0pt}
\setlength{\parskip}{5pt}
\setlength{\emergencystretch}{3em}
\setlist[itemize]{leftmargin=*,itemsep=3pt,topsep=3pt}
\setlist[enumerate]{leftmargin=*,itemsep=4pt,topsep=3pt}
\newcolumntype{L}[1]{>{\raggedright\arraybackslash}p{#1}}
\newcommand{\finding}[3]{\subsection{#1: #3}\label{rev:#1}\textbf{Priority: #2.}}
\newcommand{\file}[1]{\path{#1}}
\newcommand{\evidence}{\par\textbf{Evidence and assessment.} }
\newcommand{\action}{\par\textbf{Required revision.} }
\newcommand{\acceptance}{\par\textbf{Completion criterion.} }
\title{Independent Academic Review and Revision Recommendations\\
\large Target: IEEE Transactions on Information Forensics and Security}
\author{Manuscript under review:\\
\normalsize Allocation or Direction? A Matched-Distortion Audit\\
\normalsize of Visual Location-Privacy Mechanisms}
\date{11 September 2026}
\begin{document}
\sloppy
\maketitle

\section{Review basis, scope, and overall recommendation}

This review was written from the current submission package alone --- the
compiled \file{main.pdf} (13 pages), \file{supplementary.pdf} (6 pages) and
\file{titlepage.pdf}, read against their sources, together with the released
per-query exports, the bundled reproducibility artifact and the repository
claim auditor, both of which I re-executed. No previous review, progress
document or completion claim was used as input. Every quantitative statement
below was recomputed from the released rows with an independent
implementation --- my own paired bootstrap and signed-rank code, not the
repository's --- so that a value agreeing here agrees across two toolchains;
\S\ref{sec:verification} lists each check so the authors can repeat it. The
venue requirements in \S\ref{sec:venue} were re-verified against the IEEE
Signal Processing Society author guidance and the TIFS deep-learning
submission guidance current at the date of this review rather than assumed.

\textbf{Overall recommendation: major revision.}

Most of what I checked held, and it held exactly. Table~I reproduces cell for
cell from \file{tifs_d6}; the primary placement null reproduces under both
units of inference with every verdict unchanged; the thirteen-transform
preprocessing table reproduces in all $13\times2\times6$ cells; the four
purification runs, the gallery sweep, the clip-pooling study, the KITTI-360
replication, the factorial decomposition, the utility frontier and the
solved-map arms all reproduce to the printed precision. The protocol is a real
contribution, the negative results travel with the qualifications that bound
them, and the mechanism's own defect --- a released support map that is
numerically constant --- is disclosed rather than buried. The findings of
earlier rounds that I could re-check are closed: the white-box
``fragile to every operator'' sentence is now correct and exception-named, the
clean baseline is no longer confused with the isotropic control, the
preprocessing table is registered with the auditor, the margin oracle is
reported on the primary benchmark, the primary placement null now has a table,
the equivalence tally names its self-comparison, the two interval-unit
sentences agree, a CLIP-class attacker and a gallery sweep exist, and every
dataset and named method the paper uses is cited.

Four things nevertheless keep this below the line, and three of them need no
new computation.

\begin{itemize}
\item The paper's central surviving claim --- that a solved allocation map
buys nothing when the attacker is withheld --- is contradicted by an
experiment the authors ran, released and did not report. On CLIP ViT-L/14 the
\emph{same} surrogate-solved map that buys $-0.006$ on ResNet18 buys $-0.0508$
(place-clustered $[-0.082,-0.021]$, $p=1.6\times10^{-3}$), which is 59\% of
what the direction axis buys on that attacker (\S\ref{rev:R1}).
\item The four solved-map numbers are reported under two different reference
pairings within one paragraph, and Fig.~2 plots the pairing the same section
says must not be used. On that pairing the surrogate-solved map on ResNet18
has a place-clustered interval excluding zero, so the paper's own headline
figure disagrees with the sentence beside it (\S\ref{rev:R2}).
\item Neither optimisation arm has converged. Doubling the step budget roughly
doubles the attacker-given effect and moves the surrogate-solved arm from
$-0.0058$ to $-0.0200$ on ResNet18; the doubling is reported only where it
supports the thesis (\S\ref{rev:R3}).
\item A referee who clones the repository cannot verify 74 of the 827
registered claims, among them every cell of the factorial decomposition that
the paper lists as its fourth contribution. The rows are in the bundle, under
other names, so this is a wiring defect rather than a missing measurement ---
but TIFS states that a negative-results paper contradicting published work is
held to \emph{a higher standard of reproducibility}, and this is the paper's
own audit apparatus (\S\ref{rev:R4}).
\end{itemize}

None of the four is fatal to the contribution. \S\ref{rev:R1} in particular
makes the paper's real result sharper than the one currently claimed: the
axis asymmetry is not universal, it is attacker-dependent, and the paper has
the data to say which attackers it holds for. But the abstract, the
conclusion and one figure currently assert something the release refutes, and
that has to be fixed before this can go to a TIFS referee.

\subsection*{Multi-dimensional assessment}

\begin{center}
\begin{tabular}{@{}lcL{0.52\linewidth}@{}}
\toprule
\textbf{Dimension} & \textbf{Score} & \textbf{Basis} \\
\midrule
Originality and significance & 78/100 & The matched-delivered-distortion,
one-axis-at-a-time audit is a genuine and transferable methodological
contribution; the object audited is the authors' own testbed, which bounds the
scope of the negative result. \\
\addlinespace
Methodological rigour & 70/100 & Pre-registered margin, two units of
inference, energy-conservation gates, per-frame bisection to delivered MSE,
adaptive and purifying attackers --- against which stand
\S\ref{rev:R1}--\S\ref{rev:R3}. \\
\addlinespace
Literature coverage & 82/100 & 66 entries, 15 from 2023 or later, every
dataset and named method cited, the closest published methods adapted and run
rather than only described. \\
\addlinespace
Writing quality & 76/100 & Dense, precise and unusually honest about its own
limits; at 13 of 13 pages with no margin, and several load-bearing sentences
carry numbers from unstated variants. \\
\addlinespace
Reproducibility & 74/100 & Two checkers, 827 and 228 claims, a hash manifest
over 743 files, full software and hardware provenance --- reduced by
\S\ref{rev:R4} and by the deep-learning checklist gaps of \S\ref{rev:R5}. \\
\midrule
\textbf{Overall} & \textbf{76/100} & Major revision. \\
\bottomrule
\end{tabular}
\end{center}

\subsection*{Summary of findings}

\begin{longtable}{@{}L{0.05\linewidth}L{0.11\linewidth}L{0.60\linewidth}L{0.16\linewidth}@{}}
\toprule
\textbf{ID} & \textbf{Priority} & \textbf{Finding} & \textbf{New runs?} \\
\midrule
\endfirsthead
\toprule
\textbf{ID} & \textbf{Priority} & \textbf{Finding} & \textbf{New runs?} \\
\midrule
\endhead
R1 & Major & ``Given only surrogates a solved map buys nothing'' is false on
the third attacker the authors ran. The CLIP allocation arm is released,
protocol-identical, significant, and reported nowhere. & No \\
\addlinespace
R2 & Major & The solved-map arms are printed under two reference pairings in
one paragraph, and Fig.~2 uses the one the text rejects --- on which the
surrogate-solved interval excludes zero. & No \\
\addlinespace
R3 & Major & Doubling the search doubles the effect on every solved arm,
including the surrogate one the null rests on. The doubling is reported only
for the arms where it helps. & No \\
\addlinespace
R4 & Major & A clean clone verifies 753 of 827 claims, not 827. Five export
trees live only in a gitignored directory, though identical copies sit in the
bundle. The bundled verifier labels non-manuscript values ``paper='', and the
README's counts are stale. & No \\
\addlinespace
R5 & Moderate & Against the TIFS deep-learning checklist, three of the four
trained networks in the paper are specified only in a document that is not
part of the submission, and several required items appear nowhere. & No \\
\addlinespace
R6 & Moderate & The Patch-NetVLAD preprocessing range names the wrong upper
endpoint: 4-bit quantisation is $-0.1375$ against the printed
$-0.1350$ (median), and ``hardening roughly doubles each'' fails on two
cells. & No \\
\addlinespace
R7 & Moderate & Three Wilcoxon conventions are in use across the paper's own
tools; the same cell reads $0.38$ or $0.18$ depending on which. Intervals and
$p$-values in three tables are registered with neither checker. & No \\
\addlinespace
R8 & Moderate & The concentration argument quotes the most favourable of two
available numbers, and the literal $0.649$ denotes two different placements in
two sections. & No \\
\addlinespace
R9 & Moderate & All three hardcoded cross-document references in the
supplement are stale; two of them point at ``Table~I'' for a table the
13-page compression removed. & No \\
\addlinespace
R10 & Moderate & Responsible use and the dataset-availability explanation that
TIFS asks a negative-results submission to give its AE are supplement-only. & No \\
\addlinespace
R11 & Moderate & Nine tables the manuscript argues from are in a 15-page third
document that is not submitted; seventeen pointers out of the main body carry
no number. & No \\
\addlinespace
R12 & Minor & Packaging and mechanics: a stale \file{main.pdf}, zero margin on
all three ceilings, no biographies, six font-shape warnings, a seed-dependent
printed range. & No \\
\bottomrule
\end{longtable}

\section{Findings}

\finding{R1}{Major}{The paper's central surviving claim is contradicted by an
experiment the authors ran, released, and did not report}

\evidence
This is the objection a TIFS referee will raise first, and the evidence for it
is inside the submission's own release.

The abstract states: ``Solving for the map rather than prescribing one makes
this a fact about the rules, not the axis: given the attacker a solved map
beats uniform ($-0.039$ and $-0.094$ Top-1), \emph{and given only surrogates
it buys nothing} ($-0.006$ and $+0.006$).'' \S III-C repeats it --- ``Take the
attacker away and the axis closes'' --- and the conclusion repeats it again:
``with the attacker in hand a solved map does beat uniform \dots\ and without
it the same solver buys nothing at all.'' That asymmetry is what the paper's
thesis now rests on, and it is stated as a property of the allocation axis.

The file \file{src/exports/r10_clip/r10_clip_alloc.csv} says otherwise. It
holds 12{,}000 rows: 400 queries, three seeds, delivered MSE exactly $15.68$
on every row, energy gate \texttt{weight\_mean\_square} $=1.000000$
throughout, \texttt{optimised\_against} $=$
\texttt{expectation:resnet50+vgg16+cosplace} at 20 steps --- the same
manifest, the same budget, the same surrogate ensemble, the same step count
and the same held-out-field (\texttt{crossdraw}) scoring as the ResNet18 arm,
evaluated against CLIP ViT-L/14. Recomputing it against its own uniform arm:

\begin{center}
\footnotesize
\begin{tabular}{lccccc}
\toprule
Arm on CLIP ViT-L/14 & Top-1 & $\Delta$ & query 95\% CI & place 95\% CI & $p$ \\
\midrule
uniform (ref.) & $0.3175$ & --- & --- & --- & --- \\
solved, surrogates only & $0.2667$ & $-0.0508$ & $[-0.082,-0.022]$ & $[-0.082,-0.021]$ & $1.6\times10^{-3}$ \\
solved, surrogates, $2\times$ steps & $0.2608$ & $-0.0567$ & $[-0.089,-0.025]$ & $[-0.090,-0.024]$ & $1.0\times10^{-3}$ \\
solved, given the attacker & $0.1833$ & $-0.1342$ & $[-0.173,-0.097]$ & $[-0.174,-0.095]$ & $4.0\times10^{-10}$ \\
\midrule
direction, 3 surrogates (reported) & $0.2142$ & $-0.0867$ & $[-0.122,-0.052]$ & $[-0.126,-0.049]$ & $2.6\times10^{-6}$ \\
\bottomrule
\end{tabular}
\end{center}

Four things follow, and each of them is worse for the claim than the last.

\textbf{(a) The claim is false as stated.} On CLIP the surrogate-solved map
buys $-0.0508$ with a place-clustered interval well clear of zero. That is not
``nothing''; it is five times the $\pm0.01$ margin the paper fixed before
testing, and it is significant under the unit of inference the protocol
prescribes.

\textbf{(b) It is the same map.} The surrogate-solved arm optimises against
\texttt{resnet50+vgg16+cosplace} and never sees the evaluation attacker, so
the object released against CLIP is the same object released against ResNet18
--- confirmed by the objective trace, which falls $2.841\to2.465$ in both
files, identically. The only thing that changes between ``allocation buys
nothing'' and ``allocation buys $-0.051$'' is which held-out attacker reads
the frame. That makes the result a fact about attackers, not about the axis,
and the manuscript currently asserts the opposite.

\textbf{(c) The asymmetry the paper is built on largely collapses on this
attacker.} Direction without the attacker buys $-0.0867$ on CLIP; allocation
without the attacker buys $-0.0508$, or 59\% of it. On ResNet18 the same ratio
is 3.5\%. The sentence ``Direction differs exactly there'' is a statement
about two attackers generalised to the axis.

\textbf{(d) Nothing checks it.} The auditor registers
\file{r10_clip/r10_clip_dir.csv} and never loads
\file{r10_clip_alloc.csv}; the bundled verifier does not reach it either. A
released, protocol-identical, significant arm that refutes the abstract is
registered with neither checker, in a paper whose contribution is auditing.

I want to be explicit that this is not a request to retract the result. The
finding the data support is better than the one printed, because it is
conditional and testable: \emph{both} axes carry leverage that ordinarily
needs the attacker, direction keeps most of its when that access is withdrawn
on every attacker tested, and allocation keeps a substantial share of its own
on an attacker whose family shares nothing with the surrogates. That is a
more interesting sentence than ``allocation keeps none'', and the paper has
already paid for the experiment.

\action Report the CLIP allocation arm. Add it to \S III-C's
``Solving for the map'' paragraph and to Fig.~2's left panel, with both
intervals, alongside the two attackers already there. Rewrite the abstract's
clause and the conclusion's so that the scope is in the sentence that makes
the claim: what the evidence supports is that a solved map without the
attacker buys nothing detectable on ResNet18 and MixVPR and buys $-0.051$ on
CLIP ViT-L/14, so the direction axis is the only one that transfers on every
attacker tested, and the allocation axis transfers on some. Register both CLIP
allocation arms with the repository auditor. If the authors believe the CLIP
cell is anomalous --- it is the one attacker whose clean accuracy sits far
above ResNet18's and whose family appears in no surrogate --- say so in the
text and give the reason; do not leave the run unreported.

\acceptance The manuscript reports the surrogate-solved allocation arm on
CLIP ViT-L/14 with its place-clustered interval; no sentence in the abstract,
\S III-C or the conclusion asserts that a surrogate-solved map buys nothing
without naming the attackers that holds for; both CLIP allocation conditions
are registered claims.

\finding{R2}{Major}{Four solved-map numbers are printed under two different
reference pairings in one paragraph, and the figure uses the pairing the same
section says must not be used}

\evidence
\S III-C is explicit about how a solved map must be scored: the optimiser
redraws ``the noise field at each step so the map is a function of the
distribution rather than of any realisation'', and the arm is scored ``on a
field the optimiser never saw''. The paragraph headed \emph{What a placement
map must not see} makes the point load-bearing: a map solved against the exact
field it will be released with collapses retrieval to Top-1 $0.0025$ and
$0.0083$ and returns $+0.0025$ and $-0.0050$ on an independent field, ``so an
audit of a selective mechanism must therefore say whether its map may depend
on the perturbation it weights''. I agree, and I verified every number in
that paragraph.

The paper then does not hold to it. The released files carry both a
same-field and a held-out-field arm and both a same-field and a held-out-field
uniform reference, and the manuscript draws from both without saying which:

\begin{center}
\footnotesize
\begin{tabular}{lcc}
\toprule
Quantity ($\Delta$, place-clustered 95\% CI) & Held-out field (\texttt{crossdraw}) & Same field \\
\midrule
solved given attacker, ResNet18 & $-0.0392$ $[-0.068,-0.012]$ & $-0.0550$ $[-0.084,-0.028]$ \\
solved given attacker, MixVPR & $-0.0942$ $[-0.124,-0.065]$ & $-0.1050$ $[-0.136,-0.075]$ \\
solved from surrogates, ResNet18 & $-0.0058$ $[-0.024,+0.012]$ & $\mathbf{-0.0167}$ $\mathbf{[-0.033,-0.001]}$ \\
solved from surrogates, MixVPR & $+0.0058$ $[-0.003,+0.017]$ & $-0.0075$ $[-0.017,+0.002]$ \\
\bottomrule
\end{tabular}
\end{center}

\textbf{(a) Both appear in one paragraph.} The paragraph opens with the
held-out-field values, $-0.0392$ and $-0.0942$. Two sentences later it says
``the score-gradient placement \dots\ buys $-0.0008$ where solving that
objective buys $-0.0550$'' --- the same-field value for the same arm. Two
sentences after that, the Top-5/Top-10 sentence quotes ``$-0.078$ and $-0.071$
against $-0.055$'' and ``$-0.074$ and $-0.063$ against $-0.105$'', which I
recomputed and which are same-field throughout. A reader comparing $-0.0392$
with $-0.0550$ in the same paragraph has no way to tell they are two scorings
of one arm rather than two arms.

\textbf{(b) The figure uses the rejected pairing, and it changes a verdict.}
\file{paper/generated/fig_axes_values.tex} records exactly what Fig.~2 plots:
\texttt{solved, surrogates (ResNet18)} $=-0.0167$ $[-0.033,-0.001]$. That
interval \emph{excludes zero}. So the paper's headline figure --- the one
whose whole purpose is to put allocation and direction on one axis --- shows
the surrogate-solved allocation arm separating from zero on the weak attacker,
while the sentence three lines above it says ``no benefit detected on
either'' and the abstract says ``it buys nothing''. All four figure values are
the same-field variant: $-0.0550$, $-0.1050$, $-0.0167$, $-0.0075$.

\textbf{(c) The apparatus cannot catch this, and the file that was built to
catch it is not wired in.} \file{fig_axes_values.tex} opens with the comment
``the values plotted in the two-axis figure, so the claim auditor can verify a
figure the way it verifies a table''. The auditor contains no reference to
that file: it is not among the registry's sources and its 24 values are
unregistered. That is why the figure could drift to a different pairing from
the text without either checker noticing --- and it is the general limitation
worth stating, that a registry which verifies each number against its own
recomputation cannot detect two correctly recomputed numbers being presented
as the same quantity.

Under the held-out-field pairing the Top-5/Top-10 ordering also changes: on
ResNet18 it is $-0.0700$ at Top-5 and $-0.0758$ at Top-10 against $-0.0392$ at
Top-1, so Top-10 falls further than Top-5, which the printed pair
($-0.078$, $-0.071$) reverses. The qualitative claim survives either way; the
printed numbers do not belong to the scoring the section prescribes.

\action Choose the held-out-field pairing --- the paper's own argument
requires it --- and use it for every solved-map value in the text, in Fig.~2
and in the generated figure-value file. Where a same-field number is
deliberately kept (the $0.0025$/$0.0083$ collapse, which is the point of the
``must not see'' paragraph), label it as same-field in the sentence that
carries it. Regenerate Fig.~2 from the held-out-field values and re-check that
no plotted interval contradicts the text. Add
\file{fig_axes_values.tex} to the auditor's source list so the figure is
registered, which is what its own header promises.

\acceptance Every solved-map value in the manuscript and in Fig.~2 is drawn
from the held-out-field pairing, or is explicitly labelled otherwise; Fig.~2's
plotted values are registered claims; no plotted interval contradicts the
text beside it.

\finding{R3}{Major}{Neither optimisation arm has converged, and the doubled
budget is reported only for the arms where it supports the thesis}

\evidence
The paper answers the matched-effort objection by giving the allocation axis
``the same objective, surrogates, energy gate, delivered distortion and step
budget as the directional arm''. The released files also contain a $2\times$
budget for every solved arm, and the manuscript quotes it once: ``doubling the
search doubles both''. It does. It also doubles the arm the null rests on,
and that is not quoted.

\begin{center}
\footnotesize
\begin{tabular}{lcccc}
\toprule
Solved arm (held-out field) & 20 steps & 40 steps & place 95\% CI at 40 & $p$ at 40 \\
\midrule
given attacker, ResNet18 & $-0.0392$ & $-0.0833$ & $[-0.121,-0.046]$ & $<10^{-4}$ \\
given attacker, MixVPR & $-0.0942$ & $-0.1942$ & $[-0.234,-0.154]$ & $<10^{-15}$ \\
\emph{surrogates only, ResNet18} & $-0.0058$ & $\mathbf{-0.0200}$ & $[-0.043,+0.002]$ & $\mathbf{0.040}$ \\
\emph{surrogates only, MixVPR} & $+0.0058$ & $-0.0108$ & $[-0.026,+0.004]$ & $0.115$ \\
surrogates only, CLIP ViT-L/14 & $-0.0508$ & $-0.0567$ & $[-0.090,-0.024]$ & $1.0\times10^{-3}$ \\
\bottomrule
\end{tabular}
\end{center}

Three consequences.

\textbf{(a) The null is budget-conditional and the manuscript states it as
unconditional.} At 40 steps the surrogate-solved point estimate on ResNet18 is
$-0.0200$ --- twice the pre-registered $\pm0.01$ margin and nominally
significant by signed-rank --- and the MixVPR estimate has changed sign from
$+0.0058$ to $-0.0108$. Under the paper's own interval rule both are still
\emph{no benefit detected}, because both intervals cover zero, and I am not
claiming the null is broken at 40 steps. I am saying the trend is monotone in
the budget on all five arms and the manuscript reports it on two.

\textbf{(b) The arms are not comparable in the way ``matched budget''
implies.} The direction arm at 20 steps is at its floor: white-box Top-1
$0.0058$ on ResNet18 and $0.0008$ on MixVPR, so doubling it cannot double
anything. The allocation arm at 20 steps is on a straight line ---
$-0.0392\to-0.0833$ is $2.1\times$, $-0.0942\to-0.1942$ is $2.1\times$ --- and
its concentration is still rising, top-decile share $0.555\to0.772$ on
ResNet18 and $0.547\to0.746$ for the surrogate arm. Matching step counts
between an arm at its ceiling and an arm in its linear regime is a matched
budget in name only, and a referee will say so.

\textbf{(c) The selective quotation is the part that will be read as
unfair.} The doubling is quoted in the one sentence where it strengthens the
claim that allocation is not inert, and omitted from the sentence four lines
later where it would weaken the claim that allocation reaches nothing without
the attacker. The data for both are in the same two files.

\action Report the $2\times$ arm for every solved condition, in one small
table rather than in a clause --- five rows, two budgets, both intervals. Then
state the null at the budget it was measured at: at a matched 20-step budget a
surrogate-solved map buys nothing detectable on the two attackers tested, and
the estimate grows with the budget on all of them, so the comparison bounds
what allocation reaches at this search budget rather than what the axis can
reach. If a convergence curve is affordable --- the arms are cheap, the paper
already ran two points --- three or four step counts would settle it outright
and would be the strongest version of the experiment \S III-C is trying to be.

\acceptance The manuscript reports both step budgets for every solved arm;
no sentence states the surrogate-solved null without naming the search budget
it holds at.

\finding{R4}{Major}{A clean clone verifies 753 of 827 claims, and the 74 it
cannot reach include every cell of the paper's fourth contribution}

\evidence
The reproducibility apparatus is the strongest I have seen attached to a
submission of this kind, and it deserves saying first. The bundled manifest
verifies over 743 files; the one-command path completes in well under a minute
and reports 228 claims with 0 mismatches; the repository auditor reports 827
claims with 0 mismatched, 0 unverifiable and 0 no longer present in the
source. I ran both.

But the auditor searches \file{src/outputs/} before \file{src/exports/}, and
\file{src/outputs/} is gitignored --- it exists only on the machine that
produced it. Re-running the auditor with that directory absent, which is what
a referee who clones the repository gets:

\begin{center}
\footnotesize
\begin{tabular}{llcl}
\toprule
Missing tree & Claims lost & Where it lives instead & What it backs \\
\midrule
\texttt{tifs\_a3} & 32 & \texttt{results/factorial\_mse15} & Table~S7, operating point \\
\texttt{tifs\_a3hi} & 26 & \texttt{results/factorial\_mse241} & Table~S7, large budget \\
\texttt{tifs\_a4} & 8 & \texttt{results/maskguided\_cam25} & mask-guided comparison \\
\texttt{tifs\_a7\_o2n8} & 4 & \texttt{results/crosstime\_old2new} & cross-time subtask \\
\texttt{tifs\_a7\_n2o8} & 4 & \texttt{results/crosstime\_new2old} & cross-time subtask \\
\midrule
\multicolumn{2}{l}{\textbf{827 claims: 753 verified, 74 unverifiable}} & & \\
\bottomrule
\end{tabular}
\end{center}

The 58 claims behind Table~S7 are the whole of the factorial decomposition ---
``\emph{A decomposition} derives every control from the same optimised
perturbation and isolates which factor carries the effect'', the paper's
fourth listed contribution, and the evidence for the sentence that signs carry
the effect and magnitudes do not. The 8 behind the mask-guided table are the
paper's only head-to-head against published work.

This is a wiring defect, not a missing measurement, and it is one line to fix:
I checked the five bundle trees against the five gitignored ones and they are
identical in file count and row count ($6/19{,}200$, $6/19{,}200$, $6/7{,}200$,
$6/4{,}800$, $6/4{,}800$). Every row a referee needs is already committed,
under a reader-facing name, in \file{src/artifact/results/}. The auditor
simply does not look there.

Two further defects in the same apparatus.

\textbf{The bundled verifier prints values as ``paper='' that are not in the
paper.} Its block headed \emph{Gallery-free direction transfer (headline
table)} reports \texttt{paper=0.1508} for one surrogate, \texttt{0.0358} for
three and \texttt{0.0033} for white box on ResNet18, and
\texttt{0.7550/0.7517/0.7483/0.7342} on MixVPR. Table~I prints
\texttt{0.1608}, \texttt{0.0317}, \texttt{0.0058} and
\texttt{0.7542/0.7542/0.7467/0.7317}. Both sets are correct recomputations ---
the verifier reads
\file{results/direction_transfer_galleryfree/} and Table~I is generated from
\file{tifs_d6}, a separate execution --- but a referee running the one-command
check and comparing it to Table~I finds nine disagreements of up to $0.010$
Top-1 labelled as the headline table. That is the single most damaging thing
an auditability artifact can do.

\textbf{The README's counts are stale.} It states the auditor recomputes
\textbf{780} claims; it recomputes 827. It quotes a transcript reading
``260 distinct three- and four-decimal literals, 231 registered, 29 not'';
re-running \texttt{-{}-coverage} today prints ``268 \dots\ 236 registered, 32
not''. The README is right that the number should be a transcript rather than
a hand count; the transcript needs regenerating with the document.

Finally, the venue context. The TIFS guidance for deep-learning submissions
states: ``While we encourage the publication of negative results contradicting
previously published research, in these cases a higher standard of
reproducibility applies. Barring exceptional circumstances, data sets should be
made available for the review process (if not, this should be explained to the
AE).'' This paper is squarely in that category --- it runs adaptations of two
published mask-guided formulations and reports that neither survives its
comparison --- so the bar it is held to is the higher one, and 74 unverifiable
claims out of 827 on a clean clone is measured against that bar rather than
against the ordinary one.

\action Add \file{src/artifact/results/} as a third root for the auditor, or
copy the five trees into \file{src/exports/} under their bundle names, and
state in the README that the auditor's count is achieved from a clean clone
--- then paste the clean-clone transcript rather than the developer-machine
one. Relabel the bundled verifier's gallery-free block so it does not claim to
check the headline table: either point it at \file{tifs_d6} so it checks
Table~I's actual values, or rename it (``an independent re-execution of the
same conditions'') and print both sets with a one-line note on why they
differ, which is itself a useful reproducibility statement about optimiser
non-determinism. Regenerate the README's two transcripts. Add one paragraph,
addressed to the AE, explaining that MSLS and KITTI-360 imagery is not
redistributed under the datasets' own terms, that the released manifests carry
identifiers sufficient to rebuild every split from the public releases, and
that every reported value is recomputable from the released per-query rows
without the imagery.

\acceptance Running the auditor on a fresh clone reports 0 unverifiable
claims; no block of the bundled verifier labels a non-manuscript value
``paper=''; the README's counts and coverage transcript match what the tools
print today; the submission carries an explicit AE-facing statement on dataset
availability.

\finding{R5}{Moderate}{Against the TIFS deep-learning checklist, three of the
four trained networks are specified only outside the submission}

\evidence
TIFS maintains a separate checklist for submissions whose system relies on
deep learning, and states that the details ``could be in an appendix, or as
supplementary material providing that this is available for review''. The
package does well on several of its items and not on others.

\textbf{Satisfied in the submitted package.} GPU models per experiment family
and wall-clock cost; exact software versions (Python 3.10.12, PyTorch 2.13.0,
torchvision 0.28.0, NumPy 2.2.6, SciPy 1.15.3); attacker and surrogate
provenance with trunks named; the direction optimiser in full (20 sign-gradient
steps at step size 1 under $\ell_\infty\le16$, $\pm1$ random start, two
transforms sampled per step, BPDA substitution, bisection to delivered MSE);
the DCRF constants; and the adaptive attacker's entire recipe --- frozen
trunk, last block of 8.39\,M parameters, triplet margin $0.2$, hard-negative
mining over eight candidates, Adam at $10^{-4}$, batch 16, 30 epochs, three
mutually place-disjoint splits of 102/80/200 queries, checkpoints selected on
validation MRR. That last one is a model of what the checklist asks for.

\textbf{Not in the submitted package.} The mechanism's own unary support
predictor --- the object the whole first half of the paper audits --- is
described in \file{main.tex} as ``a small encoder--decoder of 87,441
parameters \dots\ trained with binary cross-entropy''. Its topology and
optimiser (two $3\times3$ convolution blocks at 32 and 64 channels with
max-pooling, two transposed-convolution upsampling blocks, a $1\times1$
convolution to one logit, AdamW at $10^{-4}$) appear only in the extended
evidence report, which is not part of the submission
(\S\ref{rev:R11}). The purifying attacker's network --- a twelve-layer DnCNN
with 64 channels and batch normalisation, mean absolute error on random
$128\times128$ crops, Adam at $10^{-4}$, batch 16, 8{,}000 steps --- is
likewise extended-report-only, although the submitted supplement does carry
its quality gate and its place-disjointness assertion.

\textbf{Not stated anywhere.} For the unary predictor: initialisation, batch
size, number of epochs, stopping criterion, which data determined the
hyperparameters and the hyperparameter-search policy, and the train/validation
split. For the second, mask-supervised checkpoint --- the one whose map is
genuinely non-uniform and which carries half the placement study --- nothing
beyond ``trained with structural mask supervision''. All of these are named
items on the checklist.

This matters here more than it usually would, for a specific reason. The
paper's first diagnostic is that the released support map is numerically
constant, and it uses that to argue the three redistribution controls are
near-no-ops. A referee cannot tell from the submission whether the constancy
is a property of this architecture and objective, a property of a training run
that stopped early, or a property of a checkpoint selected by a criterion that
is not stated --- and the difference bears directly on how much of the
negative result generalises.

\action Move the unary predictor's and the purifier's specifications into the
submitted supplement, and add the missing checklist items for both networks
and for the mask-supervised checkpoint: initialisation, batch size, epochs,
stopping criterion, the data used to set hyperparameters, and the split. The
supplement is at its ceiling, so this is part of the space budget in
\S\ref{rev:R11}; the operator definitions of \S VII compress most easily,
since equation-level detail there is recoverable from the released code.

\acceptance Every trained network in the paper is specified to the TIFS
deep-learning checklist within the submitted package, with no item deferred to
a document that is not submitted.

\finding{R6}{Moderate}{The Patch-NetVLAD preprocessing range names the wrong
upper endpoint}

\evidence
\S III-F reads: ``On Patch-NetVLAD the direction separates from zero under
every transform unhardened, from $-0.0133$ ($\sigma=2$ blur) to $-0.1350$
(median), and hardening roughly doubles each.''

Recomputing all thirteen from
\file{src/exports/preprocessing_r11/r5_pre_pnv_plain.csv} and
\file{r5_pre_pnv_eot.csv} (400 queries, three seeds, paired against each row's
own isotropic control):

\begin{center}
\footnotesize
\begin{tabular}{lcc|lcc}
\toprule
Transform & unhard. & hard. & Transform & unhard. & hard. \\
\midrule
none & $-0.2142$ & $-0.2292$ & resize & $-0.1000$ & $-0.1642$ \\
JPEG-75 & $-0.1250$ & $-0.2267$ & blur $\sigma{=}2$ & $\mathbf{-0.0133}$ & $-0.0242$ \\
JPEG-50 & $-0.0942$ & $-0.1650$ & 4 bits/channel & $\mathbf{-0.1375}$ & $-0.1783$ \\
blur & $-0.1117$ & $-0.2067$ & random choice & $-0.0917$ & $-0.1583$ \\
denoise & $-0.0392$ & $-0.1225$ & JPEG+blur & $-0.0758$ & $-0.1383$ \\
JPEG-60 & $-0.0883$ & $-0.1933$ & median & $-0.1350$ & $-0.2217$ \\
JPEG-30 & $-0.0558$ & $-0.1200$ & & & \\
\bottomrule
\end{tabular}
\end{center}

The lower endpoint is right: $\sigma=2$ blur at $-0.0133$ is the smallest. The
upper endpoint is not. Over the twelve non-identity transforms the largest is
4-bit quantisation at $-0.1375$, not the median filter at $-0.1350$; over all
thirteen, which is how the sentence introduces them (``We ran the same
thirteen''), the largest is the untransformed control at $-0.2142$. Under
either reading the printed maximum is not the maximum. Every cell is
significant, so the claim's substance holds --- it is the range that is
misstated, and this is the same class of defect an earlier review round
corrected on the MixVPR white-box sentence.

``Hardening roughly doubles each'' is also loose in both directions: the
untransformed row moves $1.07\times$ and 4-bit $1.30\times$, while denoise
moves $3.12\times$. A ratio range would be both more accurate and shorter than
the word ``each''.

\action Correct the endpoint --- $-0.0133$ ($\sigma=2$ blur) to $-0.1375$
(4-bit quantisation) over the twelve transforms, or state the thirteen-way
range including the untransformed control --- and replace ``roughly doubles
each'' with the measured ratio range ($1.1$--$3.1\times$, median $1.7\times$).
Register both endpoints and the ratio with the auditor; these are among the
values the coverage run currently lists as unregistered.

\acceptance The Patch-NetVLAD range names the transform set it is taken over,
its endpoints are the extreme values of that set, and both are registered
claims.

\finding{R7}{Moderate}{Three signed-rank conventions are in use across the
paper's own tools, and the intervals in three tables are registered with
neither checker}

\evidence
Two of the paper's own scripts compute the same quantity differently.
\file{src/scripts/audit_claim_consistency.py} uses
\texttt{wilcoxon(diff, zero\_method="wilcox", method="approx")} --- the normal
approximation. \file{src/scripts/make_mixvpr_placement_table.py} passes only
the non-zero differences to \texttt{wilcoxon} with default arguments, which
selects the exact test when the non-zero count is small. For Table~S5's
anti-score-gradient row the two give:

\begin{center}
\footnotesize
\begin{tabular}{lcccc}
\toprule
Cell & non-zero pairs & normal approx. & exact (generator) & printed \\
\midrule
Table~S5, anti-score-gradient & 5 of 400 & $0.180$ & $0.375$ & $0.38$ \\
Table~S4, learned support & 1 of 400 & $0.317$ & $1.000$ & $0.32$ \\
Table~S5, fixed random & 11 of 400 & $0.337$ & $0.413$ & $0.36$ \\
\bottomrule
\end{tabular}
\end{center}

No verdict changes --- every one of these cells is non-significant under every
convention --- but a referee who recomputes one number with the repository's
own auditor and compares it to the table gets a different answer, in a paper
that ships an auditor precisely so that this cannot happen. And the deeper
point is that a signed-rank $p$ computed on one discordant pair out of four
hundred, or five, is not an interpretable quantity under any convention. The
paper's interval rule is the right instrument here and it already does the
work; the $p$ column adds noise and invites exactly this kind of
inconsistency.

Separately, the auditor registers Top-1 and $\Delta$ for Tables~S2, S4 and S5
but not their confidence intervals or $p$-values
(\texttt{\_msls\_placement} and \texttt{\_mixvpr\_placement} return only
\texttt{top1} and \texttt{delta}). Those columns are the ones the verdicts are
read off, and they are unchecked.

\action Fix one convention --- I recommend the exact test on the non-zero
differences, which is what the table generators already do and what is correct
at these counts --- write it into the protocol paragraph of \S III-A, and make
\file{audit_claim_consistency.py} use the same call so the two agree. Print
the number of discordant pairs beside $p$ in Tables~S2, S4 and S5, or drop the
$p$ column from rows whose discordant count is in single figures and say why.
Register the interval endpoints and $p$-values of all three tables with the
auditor.

\acceptance One signed-rank convention is stated in the protocol and used by
every script; Tables~S2, S4 and S5 carry registered intervals and $p$-values;
no $p$ is printed for a cell with a single-figure discordant count without that
count beside it.

\finding{R8}{Moderate}{The concentration argument quotes the more favourable of
two available numbers, and one literal denotes two different placements}

\evidence
\S III-C argues that what separates the solved map from the gradient-guided
rules is not concentration: ``the margin-gradient rule packs \emph{more}
weight into its top decile than the solved map does, $0.649$ against
$0.555$''. Recomputing the top-decile weight shares actually released:

\begin{center}
\footnotesize
\begin{tabular}{lcc}
\toprule
Map & ResNet18 & MixVPR \\
\midrule
margin-gradient rule (\file{margin_oracle}) & $0.6487$ & --- \\
solved, given the attacker & $\mathbf{0.5552}$ & $\mathbf{0.6215}$ \\
solved, from surrogates & $0.5466$ & $0.5529$ \\
solved, given the attacker, $2\times$ steps & $0.7723$ & $0.7999$ \\
\bottomrule
\end{tabular}
\end{center}

$0.555$ is the ResNet18 attacker-solved map and reproduces exactly. But the
sentence sits in a paragraph that reports both attackers, and on MixVPR the
solved map's share is $0.6215$ against the margin rule's $0.6487$ --- a gap of
$0.027$ rather than $0.094$, which is not enough to carry ``concentration was
never the variable''. The argument holds on the attacker whose number is
printed and is close to vacuous on the one that is not. Nothing checks either:
both literals are in the auditor's unregistered list.

Two smaller things in the same family. The literal $0.649$ is used for the
margin-gradient rule's top-decile share in \S III-C and for the centre-bias
placement's in \S III-I; one of the two is presumably a coincidence and a
reader cannot tell which, since neither is registered. And \S III-I's edge
share of $0.839$ recomputes as $0.8404$ in the operator study's rows, which is
a different sample rather than a different quantity --- worth saying which
sample each concentration statistic is measured on, since the paper uses them
to compare rules.

\action State the concentration comparison on both attackers, or restrict the
sentence to ResNet18 and say so. Register every top-decile share the
manuscript prints --- the seven rules, the three published-model priors, the
margin oracle, and all four solved arms --- with the tree and sample each is
measured on, so that two different placements cannot print the same literal
without the checker noticing.

\acceptance Every concentration statistic in the manuscript is a registered
claim naming its source tree, and the ``concentration was never the variable''
sentence names the attackers it holds for.

\finding{R9}{Moderate}{Every hardcoded cross-document reference in the
supplement is stale, and one of them survives the compression that removed the
table it names}

\evidence
The supplement is otherwise disciplined about cross-references: every pointer
within itself goes through \verb|\ref|, and those all resolve. Its three
pointers into the manuscript are written as literals, and all three are wrong.

\textbf{(a) A section number two subsections out.} \file{supplementary.tex}
line 399, in \S X's description of Table~S6: ``Reading the two white-box
columns against each other is the point: the unhardened one is what
\verb|\S| III-G quotes''. The unhardened white-box ranges are quoted in
\S III-F, \emph{Robustness to Non-Adaptive Preprocessing}. \S III-G is
\emph{What Is Actually Released}, which concerns the amplitude bound and file
serialisation and quotes no white-box transform range at all. The pointer is
attached to the table carrying the held-out-transform claim, so it is the one
a careful referee is most likely to follow.

\textbf{(b) A table number that now names a different table --- twice.} Line
31, in the supplement's opening inventory of what did not fit: ``per-rule
placement study behind Table~I''. Line 151, opening \S VI: ``The per-rule
study behind the manuscript's Table~I \dots\ is tabulated in the extended
report, whose $\Delta$ column the manuscript's own table carries.''

The manuscript typesets exactly two tables. Table~I is the direction-transfer
table and Table~II is the utility frontier; the notation table in
\file{main.tex} sits inside a \texttt{comment} environment and the mined-proxy
placement table that used to be Table~I did not survive the compression to 13
pages. So both sentences send a referee to the direction-transfer table for
the per-rule placement study, and the second additionally describes a
$\Delta$ column that the manuscript's own table does not carry for
placements. \S III-B, which is where that study is reported, now carries the
98 comparisons in prose with no table at all --- which is correct given the
page budget, but means the supplement is describing a document structure that
no longer exists.

This is the same class of defect as \S\ref{rev:R2}: a pointer that no checker
can follow, in a package whose selling point is that everything can be
checked.

\action Replace all three literals with \verb|\ref| against labels in
\file{main.tex} --- a \verb|\S\ref| for the preprocessing section, and for the
placement study either a reference to \S III-B or, if \S\ref{rev:R11} moves
the per-rule table into the supplement, a \verb|\ref| to it there. Then grep
both documents for any remaining literal ``Table~'' or ``\verb|\S| III''
followed by a number. A \verb|\label| in the manuscript and a
\verb|\ref| in the supplement will not resolve across separate compilations
without an external-reference package, so either load \texttt{xr} against
\file{main.aux} or, at minimum, add a build-time check that asserts each
literal still names what it claims.

\acceptance No document contains a hardcoded cross-document section or table
number, or each remaining literal is asserted by a build-time check.

\finding{R10}{Moderate}{Responsible use, and the dataset-availability
statement TIFS asks a negative-results submission to give its AE, are both
supplement-only}

\evidence
The paper releases a perturbation that transfers to models its optimiser never
saw, degrades a retrieval capability investigators use legitimately, and ---
by \S\ref{rev:R1} --- transfers to a foundation-model retriever as well. The
main paper gives this two clauses inside the threat-model paragraph: ``The
same retrieval localises evidence legitimately; we release code and exports,
not perturbed imagery, and report the capability with its limits.'' The
substantive treatment --- what is and is not released, why the protocol rather
than the perturbation is regarded as the contribution, and the four limits
that bound the capability --- is the supplement's final section.

The supplement's treatment is good. Its placement is wrong for this venue.
TIFS is a security journal whose readership includes the investigators whose
capability this degrades, and dual-use framing is something its reviewers look
for in the body of the paper, not in supplementary material an AE may read
last. The same applies to the dataset statement discussed in
\S\ref{rev:R4}: the TIFS deep-learning guidance asks that a negative-results
submission explain to the AE why datasets are not available, and the
manuscript's ``No image is redistributed; the manifests carry identifiers
only'' is addressed to a reader rather than framed as that explanation.

\action Promote a short subsection --- half a column is enough --- into the
main paper, placed at the end of \S III or immediately before the conclusion:
what is released and what is withheld, the four measured limits that bound the
capability, and the reason the audit protocol rather than the perturbation is
offered as the contribution. Keep the supplement's fuller version. Add the
AE-facing dataset paragraph of \S\ref{rev:R4} to the same place or to the
cover letter.

\acceptance The main paper contains a titled responsible-use passage, and the
submission carries an explicit statement of what dataset material is and is
not available for review and why.

\finding{R11}{Moderate}{Nine tables the manuscript argues from are in a
15-page third document that is not submitted, and seventeen pointers out of the
main body carry no number}

\evidence
The supplement opens by listing what it could not fit: ``Nine more sit in the
extended evidence report for space --- per-rule placement study behind
Table~I, two-city E1 table, white-box per-backbone table, placement budget
sweep, redistribution controls by backbone, mask-guided comparison, two
utility tables, adaptive attacker's per-condition breakdown''. That report is
15 pages and ships with the code. It is not part of the submission, so a TIFS
referee reviewing the package cannot see any of it.

Several of those nine are not peripheral. The per-rule placement study is the
evidence behind Table~I, the manuscript's first table. The mask-guided
comparison is the paper's only head-to-head against published methods and the
basis for ``the prior formulation does not survive the matched-distortion
comparison''. The two utility tables carry the paired per-image measurements
that the utility argument rests on. The adaptive attacker's per-condition
breakdown is the evidence for an entire subsection.

The pointers are also unnumbered. \file{main.tex} refers to ``Supplementary
Material'' ten times and to the extended report seven times, and not once by
table number, while the supplement's tables are numbered S1--S7 and its figures
S1--S2. A referee reading \S III-F's ``the white-box bound under each transform
is tabulated in the Supplementary Material'' has to search six pages for
Table~S6.

The constraint is real: SPS allows six double-column supplementary pages
without editor approval and the supplement is at exactly six. But approval is
available for asking, and this is the case for asking --- a submission whose
contribution is an audit protocol should carry its evidence inside the
package.

\action Request EiC approval for a longer supplement and move at least the
mask-guided table, the per-rule placement study behind Table~I, and the two
utility tables into it; if approval is declined, move those four and compress
\S VII's operator definitions to pay for them. Number every pointer:
``Table~S6'', ``Fig.~S1'', through \verb|\ref| against the supplement's labels
so they cannot go stale. Where a table genuinely must stay in the extended
report, say so in the sentence that cites it, with the report named as a
release artifact rather than as ``the extended report'' alone.

\acceptance A referee can read every table the manuscript argues from inside
the submitted package, or is told explicitly and by name where a table lives
and why; every cross-document pointer carries a number.

\finding{R12}{Minor}{Packaging and submission mechanics}

\evidence
Seven items, with their fixes, in decreasing order of consequence.

\begin{enumerate}
\item \textbf{The shipped \file{paper/main.pdf} is one build behind
\file{paper/main.tex}.} The file at the submission path lacks the abstract's
utility verdict (``no budget we tested is both admissible and effective'') and
the two recent references added with it; the current build sits in
\file{paper/build/main.pdf}. Both are 13 pages. Copy the current build
forward before submitting, and make that copy part of the build script's exit
path so the two cannot diverge again. This review was written against the
current build.

\item \textbf{All three ceilings are at zero margin, and every finding above
adds text.} The main paper is 13 of 13 pages at initial submission, the
abstract 249 of 250 words, the supplement 6 of 6 pages. \S\ref{rev:R1},
\S\ref{rev:R2}, \S\ref{rev:R3}, \S\ref{rev:R5}, \S\ref{rev:R10} and
\S\ref{rev:R11} all add material. Budget them now: a revised manuscript may
reach 16 double-column pages, so the space exists, but the abstract's ceiling
does not move and the CLIP scope clause of \S\ref{rev:R1} has to be paid for
--- the mechanism-diagnostic clause (``and the learned map is numerically
constant'') is the sentence whose removal costs the abstract least, being a
finding about the testbed rather than about the question.

\item \textbf{No space is reserved for author biographies.} IEEE Transactions
print them, and the 16-page revision budget includes them. Reserve the space
deliberately rather than discovering it at acceptance.

\item \textbf{Overlength charges are already incurred.} At \$220 per page
beyond the first ten published pages, 13 pages implies roughly \$660 and every
added page raises it. Worth stating in the cover letter that the length is
deliberate, and worth weighing against what \S\ref{rev:R11} moves into the
supplement.

\item \textbf{Six font-shape warnings in the supplement.} The build log
reports \texttt{Font shape `T1/ptm/m/scit' undefined} six times --- small-caps
italic Times, silently substituted. It comes from italic text inside the
\verb|\caption| of a table typeset in small caps. Harmless to the reader,
visible to a production editor; fix by removing the emphasis from the affected
captions or loading a font that provides the shape.

\item \textbf{One printed range is bootstrap-seed dependent.} \S III-A and
\S III-C give the clustered-to-query interval-width ratio for the placement
family as $0.94$--$1.10\times$. Recomputing at 10{,}000 resamples with a
different seed gives $0.89$--$1.11\times$; the earlier review round reported
$0.92$. The verdicts are unaffected --- clustering is nearly a no-op when 206
of 277 places carry one query --- but a range quoted to two decimals should
either fix and state the resampling seed or be quoted to one.

\item \textbf{One controlled-model value rounds differently between the paper
and its own checker.} \S III-D gives the correlated operator $0.843$ where the
bundled verifier recomputes $0.842$ and passes it on tolerance. Round once,
from the checker.
\end{enumerate}

\action Apply all seven. Items 1, 5, 6 and 7 are one-line changes; items 2, 3
and 4 are decisions to take before the rewrite rather than after it.

\acceptance The shipped PDF matches the current source; the revised package
states its page budget including biographies; no build log carries a
font-shape warning; every printed range is reproducible under a stated
resampling convention.

\section{What the submission package already satisfies}
\label{sec:venue}

Verified against the IEEE Signal Processing Society information-for-authors and
the TIFS deep-learning submission guidance current at the date of this review,
so the authors do not spend revision effort re-checking these. The governing
numbers are: at most 13 double-column pages at initial submission in 10\,pt,
including title, authors and contact information, abstract, text, all images,
figures, tables, appendices and proofs, and all references; up to 16 at
revision, with appendices and proofs then counted as supplemental rather than
in the page count; supplementary material recommended at no more than six
double-column pages without explicit EiC approval; abstract 150--250 words;
single-anonymized review; and mandatory overlength charges of \$220 per page
beyond the first ten published pages.

\begin{longtable}{@{}L{0.24\linewidth}L{0.19\linewidth}L{0.51\linewidth}@{}}
\toprule
\textbf{Requirement} & \textbf{Status} & \textbf{Measured} \\
\midrule
\endfirsthead
\toprule
\textbf{Requirement} & \textbf{Status} & \textbf{Measured} \\
\midrule
\endhead
Main paper length & Satisfied, no margin & 13 pages of 13, \texttt{IEEEtran}
\texttt{[10pt,journal]} \\
Abstract length & Satisfied, no margin & 249 words; no citations, footnotes or
displayed math; no undefined abbreviations \\
Supplementary length & Satisfied, no margin & 6 pages, the ceiling without EiC
approval (content: see \S\ref{rev:R11}) \\
Review model & Satisfied & Author block, affiliations and corresponding
contact disclosed, as single-anonymized review requires \\
Keywords & Satisfied & Six index terms, all drawn from the paper's own subject
matter \\
Bibliography mechanics & Satisfied & 66 entries, 65 cited, no dangling keys,
\texttt{IEEEtran} style; every dataset (MSLS, KITTI-360, VOC, COCO, ImageNet,
ADE20K) and named method (spectral-residual saliency, DnCNN-family
purification) cited \\
Bibliography currency & Satisfied & 15 entries from 2023 or later, including
2025 and 2026 work on location leakage from multimodal models and on
adversarial attacks against place recognition \\
Cross-references & Satisfied with one exception & 0 undefined references and 0
undefined citations in all three documents; one hardcoded section number, and
it is wrong (\S\ref{rev:R9}) \\
Typesetting & Satisfied & 0 overfull boxes in main, supplement and title page;
3 and 5 underfull; six font-shape warnings (\S\ref{rev:R12}) \\
Figures & Satisfied & Both main figures present and referenced; the one-axis
audit figure is now in the main body; both supplementary figures are vector
PDF (values plotted: see \S\ref{rev:R2}) \\
Statistical significance & Satisfied & Paired tests on identical queries, two
units of inference, a margin fixed before testing, Holm correction on the
declared primary family, and exploratory tests labelled as such --- what the
TIFS guidance asks when it says authors ``should demonstrate the statistical
significance of results claimed'' (convention: see \S\ref{rev:R7}) \\
Deep-learning reproducibility & Partially satisfied & GPU models, wall clock,
exact software versions, splits, seeds and the adaptive attacker's full recipe
are given; three of four trained networks are not specified inside the
submission (\S\ref{rev:R5}) \\
Reproducibility artifact & Satisfied with caveat & Manifest verifies over 743
files; one-command check passes (228 claims, 0 mismatches); repository auditor
passes on the authors' machine (827 claims, 0 mismatched); clean-clone
behaviour and labelling: see \S\ref{rev:R4} \\
Ethics and data & Satisfied with caveat & Public datasets under their own
terms; no imagery redistributed; responsible-use section specific but
supplement-only and no AE-facing availability statement
(\S\ref{rev:R10}) \\
\bottomrule
\end{longtable}

\section{What requires new computation, and what does not}

Nothing in this review requires a GPU. Every finding is either a correction to
the text, a re-reporting of rows already released, or a wiring change in the
audit tooling. Two items would be stronger with one cheap run each, and both
are marked optional.

\begin{longtable}{@{}L{0.06\linewidth}L{0.45\linewidth}L{0.43\linewidth}@{}}
\toprule
\textbf{ID} & \textbf{Work} & \textbf{Cost} \\
\midrule
\endfirsthead
\toprule
\textbf{ID} & \textbf{Work} & \textbf{Cost} \\
\midrule
\endhead
R1 & Report the CLIP allocation arm; rescope the abstract, \S III-C and the
conclusion; register both conditions & None. Rows are in
\file{src/exports/r10_clip/r10_clip_alloc.csv}. \\
\addlinespace
R2 & One reference pairing throughout; regenerate Fig.~2; register the
figure's values & None. Both pairings are in the released files. \\
\addlinespace
R3 & Report the $2\times$ arm for every solved condition; state the null at
its budget & None. \emph{Optional:} two more step counts would give a
convergence curve --- the cheapest run in the paper. \\
\addlinespace
R4 & Add the bundle's results tree as an auditor root; relabel the verifier's
gallery-free block; regenerate the README transcripts; add the AE-facing
dataset paragraph & None. The five trees are already committed. \\
\addlinespace
R5 & Move two network specifications into the supplement; add the missing
checklist items & None, but it costs page budget. \\
\addlinespace
R6 & Correct the Patch-NetVLAD range endpoint and the doubling claim & None.
The thirteen values are recomputed in \S\ref{sec:verification}. \\
\addlinespace
R7 & One signed-rank convention across all scripts; discordant counts beside
$p$; register the three tables' intervals & None. \\
\addlinespace
R8 & Report the concentration comparison on both attackers; register every
top-decile share & None. \\
\addlinespace
R9 & One cross-reference & None. \\
\addlinespace
R10 & Promote responsible use into the main paper; add the dataset statement &
None, but it costs page budget. \\
\addlinespace
R11 & Request supplement approval; move four tables in; number every pointer &
None, but it is the largest page-budget item. \\
\addlinespace
R12 & Packaging and mechanics & None. \\
\bottomrule
\end{longtable}

\section{Verification performed for this review}
\label{sec:verification}

Every quantitative statement above was recomputed from the released exports
with an independent implementation --- my own seed-averaging, paired
bootstrap (10{,}000 resamples over queries and over places) and signed-rank
code --- rather than by re-running the repository's own analysis. The checks,
so the authors can repeat them:

\begin{enumerate}
\item \textbf{Both checkers, end to end, and once more as a clean clone.} Ran
\file{src/artifact/run_verification.sh}: \file{MANIFEST.sha256} verifies over
743 files and \file{verify_claims.py} reports 228 claims with 0 mismatches.
Ran \file{src/scripts/audit_claim_consistency.py}: 827 claims, 0 mismatched, 0
unverifiable. Re-ran the auditor with its \file{src/outputs/} root redirected
to a non-existent path, which is the state of a fresh clone: 753 verified, 74
unverifiable across five trees. Confirmed the five are present under other
names in \file{src/artifact/results/} with identical file and row counts.
Source of \S\ref{rev:R4}.

\item \textbf{Table~I, re-derived.} Recomputed every row of the transfer table
from \file{tifs_d6/d6_r18_plain.csv}, \file{d6_r18_ablation.csv},
\file{d6_mix_plain.csv} and \file{d6_mix_ablation.csv} (400 queries, three
seeds): $0.1967/0.1608/0.1300/0.0317/0.0058$ and
$0.7800/0.7542/0.7542/0.7467/0.7317/0.0008$, with every $\Delta$, interval and
$p$ to the printed precision. Then recomputed the bundle's gallery-free tree
and obtained $0.1508/0.1325/0.0358/0.0033$ --- a different execution, which is
the discrepancy in \S\ref{rev:R4}.

\item \textbf{The primary placement null, both units.} Recomputed all seven
rules from \file{icme2027_placement_msls/final/per_query.csv} against the
uniform arm, with places taken from \file{tifs_d6} (400 queries, 277 places,
206 singletons). Uniform $0.1950$; centre bias largest at $+0.0092$; score
gradient the only negative at $-0.0008$; every Top-1, $\Delta$, query interval,
place interval and verdict in Table~S4 reproduces; the widest endpoint is
$+0.0274$, matching the manuscript's $\pm0.027$; clustering changes widths by
$0.89$--$1.11\times$ and no verdict.

\item \textbf{The strong attacker's placement table.} Recomputed all seven
rows of Table~S5 from \file{placement_mixvpr_rows/per_query.csv}: uniform
$0.7800$, score-gradient $-0.0117$, every other cell to the printed precision.
Source of the $p$-value convention finding in \S\ref{rev:R7}.

\item \textbf{The solved-map arms, all four pairings.} Recomputed
the six \file{optimised_allocation/} files
(12{,}000 and 7{,}200 rows each, 400 queries, three seeds, delivered MSE
$15.68$, \texttt{weight\_mean\_square} $=1$). Held-out field:
$-0.0392$, $-0.0942$, $-0.0058$, $+0.0058$. Same field: $-0.0550$, $-0.1050$,
$-0.0167$, $-0.0075$, the last two being what Fig.~2 plots. Doubled budget:
$-0.0833$, $-0.1942$, $-0.0200$, $-0.0108$. Realised-field collapse $0.0025$
and $0.0083$ returning $+0.0025$ and $-0.0050$ on an independent draw. Sources
of \S\ref{rev:R2} and \S\ref{rev:R3}.

\item \textbf{The CLIP attacker, both axes.} Recomputed
\file{r10_clip/r10_clip_dir.csv} (direction $0.3008\to0.2142$, $-0.0867$,
place-clustered $[-0.126,-0.049]$; white box $0.0000$) and
\file{r10_clip/r10_clip_alloc.csv} (uniform $0.3175$; surrogate-solved
$0.2667$, $-0.0508$, place-clustered $[-0.082,-0.021]$, $p=1.6\times10^{-3}$;
attacker-solved $-0.1342$). Confirmed the allocation file carries the same 400
query identifiers as the primary placement manifest, the same three seeds, the
same delivered MSE and the same surrogate ensemble as the ResNet18 arm, and
the identical objective trace $2.841\to2.465$. Source of \S\ref{rev:R1}.

\item \textbf{The preprocessing family, all thirteen transforms, four
attackers.} Recomputed unhardened and hardened Top-1, the paired $\Delta$
against each row's own isotropic control, and both white-box columns for
ResNet18 and MixVPR from \file{tifs_d6/*.csv}: every cell of Table~S6
reproduces, the two ``fragile'' ranges and their named exceptions are correct,
and all sixteen held-out hardened cells exclude zero. Recomputed the same
thirteen for Patch-NetVLAD and the ViT from
\file{preprocessing_r11/*.csv}: $13/13$ significant unhardened on
Patch-NetVLAD and $2/13$ on the ViT, $12/13$ hardened on the ViT with 4-bit
quantisation the exception --- all as printed, except the Patch-NetVLAD range
endpoint of \S\ref{rev:R6}.

\item \textbf{Purification on all four attackers.} Recomputed
the four \file{purification/per_query*.csv} files (200
held-out place-disjoint queries each): ResNet18 $-0.2083\to-0.0933$, the
direction returning $0.0367\to0.1683$ ($+0.1317$), hardened $-0.1400$, the
wrong-release purifier recovering $0.0550$ of a $0.2167$ protection --- a
quarter; MixVPR $-0.0417\to-0.0217$; Patch-NetVLAD
$-0.1650\to-0.0717$ with hardening holding $-0.0750$; the ViT $-0.0150$ before
and $+0.0033$ after. All reproduce.

\item \textbf{The gallery sweep.} Recomputed
\file{r10_gallery_sweep/*.csv} over four cities and four levels: controls
$0.2783\to0.2117$ and $0.8483\to0.8083$; advantages $-0.198,-0.180,-0.177,
-0.173$ and $-0.030,-0.042,-0.038,-0.038$, each level excluding zero.

\item \textbf{The utility frontier, both columns.} Recomputed the privacy
column from \file{tifs6_a8_mse*_s2/serialized_release.csv} and the mIoU column
by accumulating per-class intersections and unions over the 200-image
\file{tifs_a5} and \file{tifs6_a5_s2} trees: all twelve Top-1 cells, all
twelve mIoU cells and clean mIoU $0.6973$ pooled ($0.697352$ and $0.697324$
per seed) reproduce. Confirmed that the two-seed privacy column is not an
avoidable shortfall: the seed-1234 runs exist but use 200 queries and, at two
budgets, a different random start, so they genuinely cannot be pooled --- the
caption's explanation is correct and would be clearer if it said this.

\item \textbf{KITTI-360, the factorial, and the clip.} Recomputed all
thirteen placement rows from \file{kitti360_rows/placement.csv} (max
$|\Delta|=0.0117$, none separating under either unit), the reference rows from
\file{baseline.csv} ($0.1498$, $0.1512$, $0.0000$) and the direction row from
\file{direction.csv} ($-0.0543$, query $[-0.097,-0.013]$, clustered
$[-0.137,+0.023]$, widening $1.91\times$). Recomputed Table~S7's operating-point
block from \file{results/factorial_mse15/} --- signs shuffled $-0.006$ and
$+0.003$, magnitudes flattened $-0.105$ and $-0.025$, placed on edges $-0.058$
and $-0.007$ --- and Table~S3 from \file{clip_pooling*/per_query.csv} at
$k=7$: mean $-0.1771$, max $-0.1944$, best frame $-0.1597$ on ResNet18 and
$-0.0391$, $-0.0859$, $-0.0347$ on MixVPR.

\item \textbf{Top-$k$, and the reachable fraction.} Recomputed
$0.1967/0.3383/0.4242 \to 0.0317/0.1025/0.1450$ and
$0.7800/0.8692/0.8917\to0.7317/0.8300/0.8600$, and confirmed the surrogate
ensemble recovers 86.5\% of the white-box effect on ResNet18 and 6.2\% on
MixVPR.

\item \textbf{Bibliography, floats and build.} 66 entries, 65 cited, none
dangling, one uncited entry harmlessly present in the \texttt{.bib}; two table
environments and two figures in the manuscript; 0 overfull boxes and 3 and 5
underfull across the main and supplement logs; six \texttt{T1/ptm/m/scit} font
warnings in the supplement; \file{paper/main.pdf} one build behind
\file{paper/main.tex}.

\item \textbf{Venue requirements.} Re-read the SPS information-for-authors for
the page, abstract, supplementary and overlength rules quoted in
\S\ref{sec:venue}, and the TIFS guidelines for deep-learning submissions for
the reproducibility checklist and the negative-results clause quoted in
\S\ref{rev:R4} and \S\ref{rev:R5}.
\end{enumerate}

\section{Assessment}

This is a careful paper and most of what I checked held exactly. The protocol
is the contribution and it is a real one: holding \emph{delivered} distortion
fixed through the clamp rather than matching a nominal $\sigma$, deriving every
control from the same optimised perturbation, separating three axes instead of
comparing whole systems, solving for the variable on the axis whose rules
failed rather than resting on the failure of prescriptions, and then applying
all of it to the authors' own testbed hard enough to expose that testbed's
learned map as a constant. Very few submissions in this area would survive the
audit this one performs on itself.

What holds it below the line is concentrated in one place, and it is the same
place three times. The paper's newest and now most load-bearing result --- what
happens when the map is solved rather than prescribed --- is reported under
two different scorings in one paragraph, is plotted in the headline figure
under the scoring the same section rejects, is stated at a search budget the
released files show it has not converged at, and is contradicted outright on
the one attacker whose family appears in no surrogate, by a run the authors
performed, released and did not report. Each of those is correctable from
rows already on disk. Together they mean that the sentence the abstract and
the conclusion now lead with --- ``given only surrogates it buys nothing'' ---
is not supported by the submission's own release.

The fix makes the paper better rather than smaller. The evidence supports a
conditional claim that is more interesting than the unconditional one: at a
matched search budget, a solved allocation map recovers nothing detectable
against two held-out retrievers and a substantial fraction of the directional
effect against a third, while the direction axis transfers on every attacker
tested and does so at a budget where it is already at its floor. That is an
axis asymmetry with a stated boundary, which is what an audit paper should be
producing, and it costs the authors one paragraph and one figure panel to say.

The fourth major finding is of a different kind and I want to be fair about
it. The reproducibility apparatus here is better than the norm at this venue
by a wide margin, and the defect is that it does not quite reach the standard
it sets for itself: a clean clone is 74 claims short, the bundled verifier
labels a different execution's values as the paper's, and the figure whose
values were generated specifically so a checker could verify them is not wired
to any checker. In a paper about auditing, held by its venue to a higher
reproducibility standard because it reports negative results against published
work, those are the findings a referee will weight most --- and all three are
an afternoon's work, because in every case the data is already committed.

\end{document}
